\documentclass[12pt,a4paper,openany,oneside]{book}

% --- Paquetes ---
\usepackage[utf8]{inputenc}
\usepackage[spanish, es-noshorthands]{babel}
\usepackage{amsmath, amsfonts, amssymb}
\usepackage{graphicx}
\usepackage{geometry}
\usepackage{hyperref}
\usepackage{booktabs}
\usepackage{fancyhdr}
\usepackage{listings}
\usepackage{xcolor}
\usepackage{tikz}
\usepackage{pgfplots}
\usepackage{colortbl}
\usepackage{multirow}
\usepackage{hhline}
\usetikzlibrary{arrows.meta, positioning, shapes.geometric, calc}
\pgfplotsset{compat=1.18}

\geometry{margin=2.5cm}

% --- Colores y estilos para código Python ---
\definecolor{codegreen}{rgb}{0,0.6,0}
\definecolor{codegray}{rgb}{0.5,0.5,0.5}
\definecolor{codepurple}{rgb}{0.58,0,0.82}
\definecolor{backcolour}{rgb}{0.95,0.95,0.92}

\lstset{
    language=Python,
    backgroundcolor=\color{backcolour},   
    commentstyle=\color{codegreen},
    keywordstyle=\color{blue},
    numberstyle=\tiny\color{codegray},
    stringstyle=\color{codepurple},
    basicstyle=\ttfamily\small,
    breaklines=true,
    frame=single,
    showstringspaces=false
}

% --- Estilo de cabecera y pie ---
\pagestyle{fancy}
\fancyhf{}
\renewcommand{\headrulewidth}{0.4pt}
\renewcommand{\footrulewidth}{0.4pt}
\fancyhead[L]{\small Carlos César Sánchez Coronel}
\fancyhead[R]{\small Sesión 8: Evaluación y Validación de Modelos}
\fancyfoot[L]{\small Carlos César Sánchez Coronel}
\fancyfoot[C]{\thepage}

% --- Portada ---
\title{
    \textbf{Sesión 8} \\ 
    \Huge \textbf{EVALUACIÓN Y VALIDACIÓN DE MODELOS} \\
    \Large Generalización, sesgo-varianza y optimización de hiperparámetros
}
\author{
    \small Por \\ \vspace{0.2cm}
    \Large \textsc{Carlos César Sánchez Coronel}
}
\date{2026}

\begin{document}

\maketitle
\tableofcontents
\addcontentsline{toc}{chapter}{Índice general}

% ------------------------------------------------------------------
% CAPÍTULO 8: SESIÓN 8 - EVALUACIÓN Y VALIDACIÓN DE MODELOS
% ------------------------------------------------------------------
\chapter{Evaluación y Validación de Modelos}

En las sesiones anteriores se han construido modelos predictivos de diversa complejidad. Sin embargo, un modelo no es útil si no se sabe cómo se comportará con datos no vistos. La evaluación rigurosa y la validación son fundamentales para garantizar la generalización, evitar el sobreajuste y seleccionar los hiperparámetros óptimos. Esta sesión profundiza en el compromiso sesgo-varianza, las técnicas de validación cruzada, las curvas de aprendizaje y los métodos de búsqueda de hiperparámetros, proporcionando las herramientas necesarias para llevar un modelo de desarrollo a producción con confianza.

\section{Logro de la sesión}
Evaluar correctamente el rendimiento de los modelos, diagnosticar problemas de underfitting y overfitting, y seleccionar los mejores hiperparámetros mediante técnicas avanzadas de validación y búsqueda.

\section{Underfitting y Overfitting}

\subsection{Definición y causas}
\begin{itemize}
    \item \textbf{Underfitting}: ocurre cuando el modelo es demasiado simple para capturar la estructura subyacente de los datos. Se manifiesta en un alto error tanto en entrenamiento como en prueba. Causas: poca complejidad del modelo, características insuficientes, regularización excesiva.
    \item \textbf{Overfitting}: ocurre cuando el modelo se ajusta demasiado a los datos de entrenamiento, capturando ruido en lugar de la señal. Se manifiesta en un error muy bajo en entrenamiento pero alto en prueba. Causas: modelo demasiado complejo, pocos datos, entrenamiento excesivo, falta de regularización.
\end{itemize}

\subsection{Diagnóstico práctico}
Para detectar estos problemas se comparan las curvas de error en entrenamiento y validación a lo largo del entrenamiento o al variar la complejidad del modelo.

Por ejemplo, en un árbol de decisión, a medida que aumenta la profundidad, el error en entrenamiento disminuye monótonamente, mientras que el error en validación primero disminuye y luego aumenta (punto de sobreajuste).

\subsection{Soluciones}
\begin{table}[h]
\centering
\begin{tabular}{|l|l|}
\hline
\textbf{Problema} & \textbf{Posibles soluciones} \\
\hline
Underfitting & Aumentar complejidad del modelo, añadir características, reducir regularización, aumentar número de iteraciones. \\
Overfitting & Reducir complejidad, aumentar regularización, aumentar datos, usar validación cruzada, early stopping, dropout (en redes), simplificar modelo. \\
\hline
\end{tabular}
\caption{Soluciones típicas para underfitting y overfitting.}
\end{table}

\section{Compromiso Sesgo-Varianza (Bias-Variance Tradeoff)}

\subsection{Descomposición del error esperado}
Para un punto fijo \(x_0\), el error cuadrático esperado de una predicción \(\hat{f}(x_0)\) se puede descomponer en:
\[
\mathbb{E}\left[ (y_0 - \hat{f}(x_0))^2 \right] = \text{Sesgo}^2 + \text{Varianza} + \text{Ruido irreducible}
\]
donde:
\begin{itemize}
    \item \textbf{Sesgo}: error debido a suposiciones erróneas en el modelo (diferencia entre la predicción promedio y el valor real).
    \item \textbf{Varianza}: variabilidad de las predicciones del modelo al cambiar el conjunto de entrenamiento.
    \item \textbf{Ruido}: parte del error que no se puede reducir por ningún modelo (inherente a los datos).
\end{itemize}

\subsection{Relación con la complejidad}
\begin{itemize}
    \item Modelos simples (alta linealidad) tienen alto sesgo y baja varianza.
    \item Modelos complejos (profundos, no lineales) tienen bajo sesgo pero alta varianza.
    \item El objetivo es encontrar el punto de equilibrio que minimice el error total.
\end{itemize}

\subsection{Visualización conceptual}
\begin{center}
\begin{tikzpicture}
\begin{axis}[
    xlabel={Complejidad del modelo},
    ylabel={Error},
    xmin=0, xmax=10,
    ymin=0, ymax=10,
    legend pos=north west,
    grid=both
]
\addplot[blue, thick] coordinates {(0,8) (2,5) (4,3) (6,2.5) (8,3) (10,4)};
\addplot[red, thick] coordinates {(0,2) (2,2.5) (4,3.5) (6,5) (8,7) (10,9)};
\addplot[green, thick] coordinates {(0,8) (2,5) (4,3.5) (6,4) (8,6) (10,9)};
\legend{Sesgo\(^2\), Varianza, Error total}
\end{axis}
\end{tikzpicture}
\caption{Compromiso sesgo-varianza. El error total mínimo se alcanza en una complejidad intermedia.}
\end{center}

\subsection{Implicaciones prácticas}
\begin{itemize}
    \item Para diagnosticar si un modelo sufre de alto sesgo o alta varianza, se observan los errores de entrenamiento y validación.
    \item Alto sesgo: errores altos en ambos conjuntos.
    \item Alta varianza: error bajo en entrenamiento, alto en validación.
\end{itemize}

\section{Validación Cruzada (Cross-Validation)}

La validación cruzada es una técnica para estimar el error de generalización de un modelo utilizando múltiples particiones de los datos, evitando depender de una única división.

\subsection{k-Fold Cross-Validation}
\begin{itemize}
    \item Se divide el conjunto de datos en \(k\) pliegues (folds) de tamaño similar.
    \item Para cada iteración \(i = 1,\dots, k\):
        \begin{itemize}
            \item Entrenar con todos los pliegues excepto el \(i\)-ésimo.
            \item Evaluar en el pliegue \(i\)-ésimo.
        \end{itemize}
    \item El error final es el promedio de los errores en cada pliegue.
\end{itemize}

\textbf{Elección de \(k\)}:
\begin{itemize}
    \item \(k = 5\) o \(k = 10\) son típicos. \(k=10\) proporciona un buen equilibrio entre sesgo y varianza de la estimación.
    \item \(k = n\) (dejando uno fuera, LOOCV) tiene bajo sesgo pero alta varianza y alto costo computacional.
\end{itemize}

\subsection{Stratified k-Fold}
Cuando las clases están desbalanceadas, es importante mantener la proporción de clases en cada pliegue. El stratified k-fold asegura que cada pliegue tenga aproximadamente la misma distribución de clases que el conjunto original. Es obligatorio en clasificación con desbalance.

\subsection{Leave-One-Out Cross-Validation (LOOCV)}
Caso extremo donde \(k = n\). Se entrena con todos menos una observación y se prueba en esa única. Tiene la ventaja de usar casi todos los datos para entrenar, pero es muy costoso y tiene alta varianza (los modelos son casi idénticos).

\subsection{Validación cruzada en series temporales}
No se puede usar k-fold estándar porque viola la dependencia temporal. Se usa \textbf{Time Series Split} o \textbf{expanding window}: los datos de entrenamiento siempre son anteriores a los de validación.

\subsection{Estimación del error y su varianza}
El error de validación cruzada es un estimador insesgado del error de generalización (aunque con cierta varianza). Se puede calcular la desviación estándar de los errores de los pliegues para tener una idea de la estabilidad.

\subsection{Ejemplo numérico}
Supongamos un dataset con 100 observaciones, usamos 5-fold. Cada pliegue tiene 20 observaciones. Entrenamos 5 modelos y obtenemos errores: [0.15, 0.12, 0.18, 0.14, 0.16]. El error promedio es 0.15 y la desviación típica $\approx$ 0.02. Esto indica que el modelo tiene un error esperado alrededor de 0.15 con poca variabilidad.

\section{Curvas de Aprendizaje}

Las curvas de aprendizaje muestran la evolución del error en entrenamiento y validación a medida que se incrementa el tamaño del conjunto de entrenamiento. Son útiles para diagnosticar sesgo y varianza.

\subsection{Construcción}
Para diferentes tamaños de entrenamiento (por ejemplo, 10\%, 20\%, ..., 100\% de los datos), se entrena el modelo y se evalúa en el conjunto de entrenamiento (o una parte fija de validación). Se grafican ambas curvas.

\subsection{Interpretación}
\begin{itemize}
    \item \textbf{Alto sesgo}: las curvas de entrenamiento y validación convergen a un error alto, y el error de entrenamiento es también alto. Añadir más datos no ayuda.
    \item \textbf{Alta varianza}: hay una gran brecha entre las curvas de entrenamiento (bajo error) y validación (alto error). Añadir más datos puede reducir la brecha.
    \item \textbf{Modelo ideal}: ambas curvas convergen a un error bajo y están cerca.
\end{itemize}

\begin{center}
\begin{tikzpicture}
\begin{axis}[
    xlabel={Tamaño del conjunto de entrenamiento},
    ylabel={Error},
    xmin=0, xmax=100,
    ymin=0, ymax=1,
    legend pos=north east,
    grid=both
]
\addplot[blue, thick] coordinates {(10,0.1) (20,0.08) (40,0.07) (60,0.065) (80,0.063) (100,0.06)};
\addplot[red, thick] coordinates {(10,0.6) (20,0.4) (40,0.3) (60,0.25) (80,0.23) (100,0.22)};
\legend{Entrenamiento, Validación}
\end{axis}
\end{tikzpicture}
\caption{Curvas de aprendizaje típicas de un modelo con alta varianza (gran brecha). Añadir datos ayuda.}
\end{center}

\subsection{Uso práctico}
\begin{itemize}
    \item Si el error de validación es alto y el de entrenamiento bajo $\rightarrow$ alta varianza.
    \item Si ambos errores son altos $\rightarrow$ alto sesgo.
    \item Si ambos errores convergen a un valor bajo pero hay brecha $\rightarrow$ aún se puede mejorar con más datos.
\end{itemize}

\section{Búsqueda de Hiperparámetros}

Los hiperparámetros son parámetros que no se aprenden durante el entrenamiento (ej. \(k\) en KNN, profundidad en árboles, learning rate). La búsqueda de hiperparámetros consiste en encontrar la combinación que optimiza el rendimiento en validación.

\subsection{Grid Search}
\begin{itemize}
    \item Se define una rejilla de valores para cada hiperparámetro.
    \item Se evalúa mediante validación cruzada cada combinación.
    \item Se selecciona la combinación con mejor rendimiento.
\end{itemize}
\textbf{Ventaja}: exhaustivo. \textbf{Desventaja}: costo exponencial en número de hiperparámetros. Inviable para muchos parámetros.

Ejemplo:
\begin{lstlisting}[language=Python]
param_grid = {
    'max_depth': [3, 5, 7],
    'learning_rate': [0.01, 0.1, 0.3],
    'n_estimators': [100, 200]
}
# 3*3*2 = 18 combinaciones.
\end{lstlisting}

\subsection{Random Search}
\begin{itemize}
    \item Se define una distribución para cada hiperparámetro (continua o discreta).
    \item Se muestrean aleatoriamente combinaciones y se evalúan.
    \item Teóricamente, con \(n\) iteraciones, se explora mejor el espacio que grid search, especialmente cuando algunos hiperparámetros son más importantes que otros.
\end{itemize}
\textbf{Ventaja}: más eficiente que grid search en espacios de alta dimensionalidad. \textbf{Desventaja}: no garantiza encontrar el óptimo, pero con suficientes iteraciones se aproxima.

\subsection{Bayesian Optimization}
\begin{itemize}
    \item Construye un modelo probabilístico (usualmente Gaussian Process) de la función de rendimiento respecto a los hiperparámetros.
    \item Utiliza una función de adquisición para decidir qué combinación probar a continuación, balanceando exploración y explotación.
    \item Es más eficiente que random search porque aprende de las evaluaciones previas.
\end{itemize}
\textbf{Implementaciones}: Scikit-Optimize, Hyperopt, Optuna.

Ejemplo con Optuna:
\begin{lstlisting}[language=Python]
import optuna

def objective(trial):
    params = {
        'max_depth': trial.suggest_int('max_depth', 3, 10),
        'learning_rate': trial.suggest_float('learning_rate', 0.01, 0.3, log=True),
        'n_estimators': trial.suggest_int('n_estimators', 100, 500)
    }
    model = XGBoost(**params)
    score = cross_val_score(model, X, y, cv=5).mean()
    return score

study = optuna.create_study(direction='maximize')
study.optimize(objective, n_trials=50)
best_params = study.best_params
\end{lstlisting}

\subsection{Comparación de métodos}
\begin{table}[h]
\centering
\begin{tabular}{|l|c|c|c|}
\hline
\textbf{Método} & \textbf{Eficiencia} & \textbf{Riesgo de óptimo local} & \textbf{Facilidad de implementación} \\
\hline
Grid Search & Baja (exhaustivo) & Ninguno & Alta \\
Random Search & Media & Bajo & Alta \\
Bayesian Optimization & Alta & Bajo & Media (requiere librería) \\
\hline
\end{tabular}
\caption{Comparación de métodos de búsqueda de hiperparámetros.}
\end{table}

En la práctica, para modelos con pocos hiperparámetros, grid search es aceptable. Para modelos más complejos, random search o bayesian optimization son preferibles.

\section{Consideraciones prácticas sobre la validación}

\subsection{Data Leakage}
Uno de los errores más comunes es dejar que información del conjunto de prueba o validación influya en el entrenamiento. Ejemplos:
\begin{itemize}
    \item Escalar datos con estadísticas de todo el dataset antes de dividir.
    \item Usar validación cruzada para seleccionar características y luego evaluar en el mismo conjunto.
    \item Hacer oversampling (SMOTE) antes de dividir, mezclando información de prueba.
\end{itemize}
\textbf{Solución}: siempre dividir primero, y aplicar cualquier transformación (normalización, imputación, selección) ajustándola solo en el conjunto de entrenamiento y luego transformando validación/prueba con esos parámetros.

\subsection{Validación en problemas con grupos}
Cuando los datos tienen grupos naturales (ej. varios registros del mismo paciente), la validación cruzada debe respetar esos grupos para no tener información del mismo sujeto en entrenamiento y validación. Se usa \textbf{GroupKFold}.

\subsection{Estrategia final: entrenar con todos los datos}
Una vez seleccionados los hiperparámetros, se entrena el modelo final con todos los datos disponibles (entrenamiento + validación) y se evalúa en el conjunto de prueba reservado (hold-out). El conjunto de prueba solo se usa una vez al final.

\section{Caso integrado: Selección de hiperparámetros para un XGBoost en un problema de clasificación}

\textbf{Problema}: Predecir si un cliente comprará un producto (clasificación binaria) con 50,000 observaciones y 20 características. Las clases están balanceadas.

\textbf{Pasos}:
\begin{enumerate}
    \item Dividir en entrenamiento (80\%) y prueba (20\%).
    \item Sobre el conjunto de entrenamiento, aplicar validación cruzada estratificada de 5 folds para evaluar combinaciones.
    \item Realizar una búsqueda con Optuna (50 trials) sobre:
        \begin{itemize}
            \item \texttt{max\_depth}: entero entre 3 y 10
            \item \texttt{learning\_rate}: log-uniforme entre 0.01 y 0.3
            \item \texttt{subsample}: uniforme entre 0.6 y 1.0
            \item \texttt{colsample\_bytree}: uniforme entre 0.6 y 1.0
            \item \texttt{lambda}, \texttt{alpha}: log-uniforme entre 1e-8 y 10
            \item \texttt{n\_estimators}: entero entre 100 y 500 (con early stopping interno)
        \end{itemize}
    \item Evaluar cada combinación con la media del AUC en los 5 folds.
    \item Seleccionar la mejor combinación.
    \item Entrenar un modelo final con esos hiperparámetros usando todo el conjunto de entrenamiento.
    \item Evaluar en el conjunto de prueba (hold-out) reportando AUC, precisión, recall, etc.
\end{enumerate}

\textbf{Resultados esperados}: El modelo final tendrá un AUC en prueba cercano al estimado por validación cruzada (si no hay leakage). Las curvas de aprendizaje pueden mostrar si con más datos se podría mejorar.

\section{Métricas en validación cruzada}

Dependiendo del problema, se utiliza la métrica adecuada (RMSE, AUC, F1, etc.) y se promedia sobre los folds. Es importante reportar tanto la media como la desviación estándar para entender la estabilidad del modelo.

\textbf{Ejemplo de reporte}:
\begin{itemize}
    \item AUC: 0.845 ± 0.012
    \item Precisión: 0.78 ± 0.03
    \item Recall: 0.81 ± 0.04
\end{itemize}

\section{Conexión con el resto del curso}

La evaluación y validación son transversales a todos los modelos. Los conceptos de esta sesión se aplicarán en:
\begin{itemize}
    \item Selección de hiperparámetros en redes neuronales (sesiones futuras).
    \item Comparación de modelos en proyectos finales.
    \item Detección de overfitting en modelos complejos como Gradient Boosting.
\end{itemize}

\section{Resumen}
\begin{itemize}
    \item Underfitting y overfitting se diagnostican comparando errores de entrenamiento y validación.
    \item El compromiso sesgo-varianza guía la elección de la complejidad del modelo.
    \item La validación cruzada (k-fold, estratificada, LOOCV) estima el error de generalización de forma robusta.
    \item Las curvas de aprendizaje ayudan a decidir si conviene agregar más datos o aumentar complejidad.
    \item Grid search, random search y bayesian optimization son métodos para buscar hiperparámetros, con diferente eficiencia.
    \item Es crucial evitar el data leakage y respetar la estructura de los datos (grupos, tiempo).
\end{itemize}

\end{document}